\section{Experimental Protocol}
\label{sec:experimental_protocol}
With the formal game and its system architecture fully defined, we now turn to how experiments are \textbf{systematically designed, executed, and evaluated} using \game.  
The protocol below is informed by best practices from causal benchmarking initiatives \citep{CauseBox2021,CausalBench2024,RealCause2021} and provides a reproducible workflow that any lab can adopt without reverse-engineering hidden scripts.

\subsection{Step 1: Declare the Problem Instance}
Every study begins by instantiating a \textbf{problem instance} via Specs (Sec.~\ref{sec:specs}).  
Researchers document:
\begin{itemize}
    \item the \textbf{mission specification} (e.g.\ DAG discovery, CATE estimation, SCM reconstruction);
    \item the \textbf{ground-truth SCM} (dataset-derived, Bayesian-network-based, or physics-inspired);
    \item \textbf{initial conditions} $\initstate$ and background observations shared with the agent;
    \item \textbf{stopping conditions} $\phi$ capturing sample, round, or time budgets;
    \item one or more \textbf{evaluation metrics} $(\resultmetric, \behaviormetric)$;
    \item optional \textbf{hooks} for logging or auditing.
\end{itemize}
Because Specs are plain text artefacts, the declaration can be pre-registered alongside a paper submission or shared as supplementary material—mirroring the transparency expected in clinical or social science trials \citep{Prosperi2020CausalIAA,Athey2016TheSOA}.

\subsection{Step 2: Select and Configure Agents}
Agents are instantiated via Specs under the \texttt{Agent} contract.  
Each agent configuration explicitly captures:
\begin{itemize}
    \item random seeds and exploration strategies,
    \item algorithmic hyperparameters (e.g.\ maximum experiments per round, posterior update cadence),
    \item optional prior knowledge modules or warm-start data.
\end{itemize}
No agent may modify the problem instance or SCM at runtime; all causal interaction is mediated through DTOs. This guarantees that a comparison between, say, a Bayesian design agent and a reinforcement-learning agent is fair, because both operate under identical environmental laws.

\subsection{Step 3: Execute the Game}
Running the experiment is as simple as invoking
\[
    \texttt{Game.run()} \;\;\longrightarrow\;\; \texttt{Transcript},
\]
yet the side effects are substantial. The runtime produces a \textbf{sealed interaction log} containing:
\begin{itemize}
    \item every action issued by the agent,
    \item the associated interventional samples,
    \item budget consumption trajectories and triggered penalties,
    \item the final stop decision or forced termination event.
\end{itemize}
When multiple agents, SCMs, or random seeds must be evaluated, the \texttt{Runner} (Sec.~\ref{sec:runner}) orchestrates the grid and ensures isolation between runs.

\subsection{Step 4: Evaluate and Analyse}
Evaluation is strictly \textbf{post-hoc}.  
From the transcript and ground-truth SCM, the system computes
\[
    \big( \resultmetric(y),\ \behaviormetric(r) \big),
\]
which researchers may keep as a Pareto pair or aggregate using domain-specific preferences (lexicographic ordering, convex combinations, dominance filters). Annotated transcripts can be replayed to inspect intervention sequences, allowing diagnostic comparisons similar to those advocated in \cite{CausalDiscoveryFramework2024}.

\subsection{Step 5: Disseminate Artefacts}
The protocol encourages publishing the full experiment bundle:
\begin{itemize}
    \item Spec files for problem instances, agents, missions, and metrics;
    \item transcripts and summary statistics (accuracy vs.\ cost curves, Pareto frontiers);
    \item optional plots and dashboards generated post-hoc;
    \item verification scripts to re-run the complete study.
\end{itemize}
These artefacts enable third parties to reproduce results, stress-test agents under new budgets, or extend missions without re-implementing infrastructure—directly addressing reproducibility concerns raised by prior benchmark audits \citep{CIPCaDBench2022,EvaluationFramework2022}.

\medskip
Following this protocol transforms \game from a formal abstraction into a \textbf{repeatable scientific experimentation process}. In the next section we instantiate the workflow with concrete \textbf{mission examples} and practical lessons learned from running them at scale.
